\title{Focused Transformer: Contrastive Training for Context Scaling}

\begin{abstract}
Large language models have an exceptional capability to incorporate new information in a contextual manner. However, the full potential of such an approach is often restrained due to a limitation in the effective context length. One solution to this issue is to endow an attention layer with access to an additional context, which comprises of (key, value) pairs. Yet, as the number of documents increases, the proportion of relevant keys to irrelevant ones decreases, leading the model to focus more on the irrelevant keys. We identify a significant challenge, dubbed the \emph{distraction issue}, where keys linked to different semantic values might overlap, making them hard to distinguish. To tackle this problem, we introduce the \emph{Focused Transformer} (\textsc{FoT}), a technique that employs a training process inspired by contrastive learning. This novel approach enhances the structure of the (key, value) space, enabling an extension of the context length. Our method allows for fine-tuning pre-existing, large-scale models to lengthen their effective context. This is demonstrated by our fine-tuning of $3 B$ and $7 B$ OpenLLaMA checkpoints. The resulting models, which we name \textsc{LongLLaMA}\footnote{We release the \href{https://github.com/CStanKonrad/long_llama}{checkpoints and source code} of \textsc{LongLLaMA} \raisebox{-3pt}, see also our {colabs}.}, exhibit advancements in tasks requiring a long context. We further illustrate that our \textsc{LongLLaMA} models adeptly manage a $256 k$ context length for passkey retrieval.
\end{abstract}

\section{Introduction}
Language models have served as a catalyst for substantial advancements in several areas, including natural language processing \citep{radford2019language, brown2020language}, code generation \citep{chen2021evaluating, li2022competition}, quantitative reasoning \citep{lewkowycz2022solving} and theorem proving \citep{polu2020generative, jiang2022thor, mikula2023magnushammer}. One of the central challenges with language models is the effective incorporation of extensive new knowledge. The common practice of fine-tuning the model is not only resource-intensive and complex to manage, but it also does not always clearly indicate how to incorporate new knowledge. For example, fine-tuning on a text such as ``Alice in Wonderland'' does not equip the model to answer questions about the story itself, but rather it trains the model to predict the next token or complete masked sentences. A promising alternative – integrating the new knowledge within the context – doesn't require training but is considerably restricted by the model's effective context length. For this method to work with large knowledge databases ({like large code repositories)}, the model needs to manage a context length extending to millions of tokens.

In this research, we highlight one of the primary obstacles in augmenting the context length: as the number of documents increases, the ratio of pertinent to irrelevant tokens diminishes. The standard training procedure frequently results in overlaps between keys connected with irrelevant values and those related to relevant ones, exacerbating the model's task of differentiating between them. We term this challenge the \emph{distraction issue}.

We propose the \emph{Focused Transformer} (\textsc{FoT}), an innovative technique developed explicitly to address this issue. The Focused Transformer permits a subset of attention layers to access an additional context of (key, value) pairs through the k-nearest neighbors (kNN) algorithm, akin to the method used in \citep{wu2022memorizing}. This mechanism effectively extends the total context length. The distinctive aspect of the Focused Transformer is its training procedure, drawing from contrastive learning. This method addresses the distraction issue and facilitates larger context capacities. Specifically, during the training phase, we deliberately expose the chosen subset of attention layers to both relevant and irrelevant keys (like negative samples from unrelated documents). This strategy incentives the model to differentiate keys connected with semantically diverse values, thereby enhancing their structure.

We introduce and make available \textsc{LongLLaMA}s (), fine-tuned OpenLLaMA models with \textsc{FoT}, demonstrating that our method does not require long context during training and can be applied to existing models. Notably, \textsc{LongLLaMA}s show significant improvements on tasks necessitating long-context modeling. In particular, they can manage a $256k$ context length on the passkey retrieval task \citep{mohtashami2023landmark}.

Our research contributions are the following:

\textbf{1.}  We pinpoint the  distraction issue as a significant challenge and a primary obstacle to scaling up the context length in Transformer models, particularly in multi-document scenarios.

\textbf{2.} We develop the Focused Transformer (\textsc{FoT}), designed to alleviate the distraction issue. \textsc{FoT} includes a unique training objective that improves the (key, value) structure, enabling the use of extensive additional context and k-nearest neighbors lookup to scale the context length.

\textbf{3.} Our method is simple to implement, and it provides the benefit of {extending model context} without modifying {the} architecture, facilitated by cost-effective fine-tuning. We demonstrate this on the $3B$ and $7B$ OpenLLaMA checkpoints. The resulting models, named  \textsc{LongLLaMA}s, display enhancements on tasks that benefit from increasing the number of few-shot demonstrations in the extended context, such as TREC \citep{li-roth-2002-learning, hovy-etal-2001-toward} and WebQS \citep{berant-etal-2013-semantic}. We also prove that for passkey retrieval \cite{mohtashami2023landmark}, our \textsc{LongLLaMA} models successfully handle a $256k$ context length. 

\textbf{4.} We further scrutinize  \textsc{FoT}'s capabilities across various datasets and model sizes. We show that a \textsc{FoT} trained with a total context of $512$ tokens can extrapolate to $16$ million tokens in a benchmark dictionary lookup task. We also assess \textsc{FoT} on long-context language modeling tasks such as books (PG-19), mathematics (arXiv), code (GitHub), and formal proofs (Isabelle), where it exhibits improvements in perplexity over baselines.

\section{Related work} \label{sec:related_work}

 \paragraph{Long-context transformer architectures} A multitude of  approaches have been developed to increase the context length of transformers, mostly focusing on alleviating the quadratic complexity of the attention computation. For instance, Transformer-XL \citep{dai2019transformerxl} caches the previous context and enables the linear extension of context with the number of layers. Longformer \citep{beltagy2020longformer} employs an attention mechanism that allows tokens to attend to distant tokens sparsely, reducing the computational complexity. BigBird \citep{zaheer2020bigbird}, LongT5 \citep{guo2021longt5}, and  \citep{dao2022flash} also use sparse attention to handle long sequences. Different efficiency considerations have been studied in \citep{kaddour2023train}, showing that they lead to limited gains. Hierarchical transformers \citep{nawrot2021hierarchical, nawrot2023efficient} downsample activations in intermediate layers to reduce computation and enable longer contexts. COLT5 \citep{ainslie2023colt5} proposes conditional computation to save memory and enable larger contexts. Memorizing Transformer \citep{wu2022memorizing} uses kNN lookup to pick up the most relevant tokens, which might also be seen as a way to reduce the computational complexity of attention. Our work adheres to this approach and aims to train a key space that handles longer attention context length (e.g., by mitigating the distraction issue) and, thus, has better long-context capabilities.

\paragraph{Fine-tuning LLMs for longer {retrieval}} Prior works such as RETRO \citep{borgeaud2022retro} (RETROfitting) and Memorizing Transformer \citep{wu2022memorizing} have demonstrated a promising path for fine-tuning existing LMs to add new capabilities without the need to retrain the entire model. {In contrast to those approaches our method is not framed as a retrieval but as a way of extending the context of the model. In contrast to RETRO, we propose a single-stage method for context extension instead of a two-stage retrieve-then-embed approach. We provide a more detailed comparison with the Memorizing Transformer in Appendix \ref{sec:relation_to_mt}.} More recently, a number of works have explored fine-tuning LLaMA to extend its context length. Landmark attention \citep{mohtashami2023landmark} proposes a compression scheme of LLM's context into landmarks, increasing the context length of LLaMA-7B to $32K$. Position Interpolation (PI, \citep{chen2023extending} and \citep{superhot}) introduces a modification to the rotary positional encoding scheme that enables fine-tuning for $32K$ context. In contrast to this work, our method does not rely on positional encodings, following the findings from \citep{haviv2022transformer}. Removing positional encoding in additional context allows us to extrapolate to $256k$ tokens, although the model was only trained on sequences up to $8K$, yielding theoretically unbounded context length.

\paragraph{Zero-shot methods} {KNN-LM \citep{khandelwal2019knnlm} shows that one can improve the performance of a LLM by combining two probability distributions. One created by a pre-trained model, and one based on the similarity between the embedding of the currently processed token and the embeddings of tokens retrieved from a large database. Meanwhile, we extend the model context in a subset of attention layers, potentially allowing for reasoning within this extended context.} {Parallel Context Windows for Large Language Models \citep{ratner2023parallel} introduces a method for extending the context of language models without training. They achieve this by embedding several context windows independently in parallel and allowing only a subset of tokens to attend to all windows. On the other hand, we fine-tune existing models and allow all tokens to attend to all previous tokens but only in a subset of layers. Additionally, our method allows us to improve the structure of the key-value space of the existing models.}

\paragraph{Contrastive learning} Contrastive learning aims to learn good representations by comparing positive and negative examples. CLIP \citep{radford2021clip} and SimCLR \citep{chen2020simclr} are two popular contrastive learning methods that have achieved state-of-the-art performance in the image domain. During contrastive pre-training, negative examples are kept in the same batch to learn to distinguish them from positive examples. Scaling the batch size in contrastive learning has been demonstrated to enhance the quality of representations, as shown in \citep{gao2021scaling}. It has been suggested \citep{gao2019representation} that the embedding space in language modeling suffers from degeneracy, where embeddings are tightly packed in a narrow cone, making it difficult to distinguish between them. TRIME \citep{zhong2022training} proposes a training approach designed for training LMs with memory augmentation, which uses negatives to improve the quality of representations. The main difference between this and our approach is that we incorporate negatives into the chosen subset of attention layers instead of interpolating in the output layer {and use the standard language modeling loss}. {TRIME \citep{zhong2022training} also focuses on retrieval from large databases, whereas we focus on extending the context of the model.} ContraCLM \citep{jain2023contraclm} applies contrastive losses at both the token and sequence levels during training to promote more uniformly distributed, isotropic representations. It is shown to enhance the discrimination of representations on textual semantic similarity benchmarks. While ContraCLM focuses on improving the general expressiveness of representations, our work introduces contrastive-inspired techniques designed specifically for training the attention mechanism to handle longer context lengths. Nonetheless, exploring other contrastive learning objectives could be beneficial for further improving the key structure in future work.

\section{\textsc{FoT}: Focused Transformer} \label{section:methods}

Our method, the Focused Transformer (\textsc{FoT}), is a simple plug-and-play extension of transformer models and can be used both to train new models or fine-tune existing, possibly large, models with longer context. To this end, \textsc{FoT} uses \emph{memory attention layers} and the \emph{crossbatch} training procedure. Memory attention layers enable the model to retrieve information from the additional context at inference time, effectively extending the context. The crossbatch training procedure biases the model to learn $(key, value)$ representations, which are easy to use by a memory attention layer. See Figure~\ref{fig:crossbatch} for an overview of the \textsc{FoT} architecture and Appendix \ref{appendix:code} for pseudocode. 

\subsection{Memory attention layers}\label{subs:RetrievalAugmentation}
Memory attention layers $\mathcal{L}$ are endowed with access to an additional context during inference. Namely, each query in $\ell\in \mathcal{L}$ attends to preceding keys from the local context and the top $k$ most matching keys (i.e. having the largest inner product with the query) from memory. The memory keys are ranked by the inner product with the query and retrieved using the kNN search algorithm.  We use the exact kNN search implemented in FAISS \citep{faiss_johnson2017billionscale}. The memory is populated incrementally with $(key, value)$ pairs processed by $\ell$ beforehand. Our memory attention layer design is closely related to \citep{wu2022memorizing}, we follow most of its design choices, except for the gating, which we replace with a simpler mechanism, which turns out to be more effective in our applications. See details in Section \ref{sec:relation_to_mt} and  Appendix \ref{appendix:mem_attetnion}. 
We remove positional encodings in memory layers in all our models except \textsc{LongLLaMA}s. This allows \textsc{LongLLaMA} checkpoints to be a drop-in replacement for LLaMA checkpoints. {We treat the kNN search algorithm as an approximation of full dense attention, which opens the doors for future speed-ups.}

\subsection{Crossbatch training procedure} \label{sec:training_procedure}

Our  training procedure is a novel way of training (or fine-tuning) transformer-based architectures in order to improve the structure of the $(key, value)$ space. The main motivation is to shape this space so that a memory attention layer $\ell \in \mathcal{L}$ can easily focus on relevant information. The key idea, inspired by contrastive learning, is to expose $\ell$ to $(key, value)$ pairs from the current and previous local context of the given document (positives) and $d-1$ contexts from unrelated documents (negatives). Importantly, this is done in a differentiable way.

To achieve this, we use a data pipeline in which each element of the batch corresponds to a different document. We embed the previous ($C_{\mathrm{prev}}$) and the current ($C_{\mathrm{curr}}$) local context for each of the processed documents. The overview of our procedure can be found in Figure \ref{fig:crossbatch}. Specifically for each document $\delta$ in $C_{\mathrm{curr}}$ we create a set $\{p^\delta_i\}_{i=\{1, \ldots, d\}}$  consisting of the $(key, value)$ pairs from the previous local context of $\delta$ (positives), along with pairs from $d-1$ other contexts coming from $C_{\mathrm{prev}}$ (negatives). We also experiment with varying the number of previous contexts and negatives for different batch elements.

The operation is fully differentiable, and thus, we improve all the $(key, value)$ pairs in $p^\delta$.
Two, the procedure is easy to implement; it does not require any additional loss (i.e., uses the standard transformer training objective) and is done on the level of the data loading pipeline and a minor self-attention change. The only new hyperparameter is $d$, which prescribes the ratio of positive to negative samples. Typically, we find it beneficial to start with small $d \leq 8$ (otherwise, the model tends to ignore the previous local context) and later switch to bigger values, say $d\geq64$. Appendix~\ref{appendix:training_procedure} provides more details about the method. {Listing \ref{listing:example_cb_implemtation} outlines an  implementation of the crossbatch.}

\subsection{The distraction issue} \label{sec:distraction_issue}

In this section, we conceptualize what we call the distraction issue and hypothesize it is one of the key problems in {dealing with long multi-document contexts (like large code repositories)}. Namely, during the standard training, the model is not incentivized to distinguish the keys from different documents. We measure that the attention mass is evenly spread on the related and unrelated documents; see Figure~\ref{fig:distraction issue}. More precisely, for a document~$\delta$, let $w_{ij}$ be the softmax weights related to $p^\delta_{ij}$ constructed as described in Section \ref{sec:training_procedure}. We define the positive attention mass as $r_d:=\sum_{j} w_{1 j}/ \sum_{i=1}^d\sum_{j} w_{ij}$. We observe that $r_d \approx 1/d$, which can be interpreted as the fact that the attention is equally distracted by the positive (coming from the current document at $i=1$) and negative keys. This is an undesirable property since when scaling the memory, the attention becomes increasingly distracted. We show that the crossbatch mostly alleviates the distraction issue, resulting in a \emph{focused} attention. More information can be found in Appendix \ref{appendix:distraction_issue}. In Section \ref{sec:distractions}, we also show that the distraction issue has a harmful effect on metrics like perplexity.

\section{\textsc{LongLLaMA} \raisebox{-2pt}: extending LLaMA's context length with \textsc{FoT}}\label{sec:longLama}
One of the promises of our work is that \textsc{FoT} can be used to fine-tune already existing large models to extend their context length. In this section, we show that this is indeed the case. We use \mbox{OpenLLaMA-3B} and OpenLLaMA-7B models trained for $1T$ tokens as starting points and fine-tune them with \textsc{FoT}. We show that the resulting models, which we call \textsc{LongLLaMA}s, are capable of extrapolating beyond their training context length (even up to $256K$) and retain the performance on short-context tasks. We release the {inference} code on GitHub: \href{https://github.com/CStanKonrad/long_llama}{\texttt{https://github.com/CStanKonrad/long\_llama}} and the \textsc{LongLLaMA}{-3B} checkpoint on Hugging Face: \href{https://huggingface.co/syzymon/long_llama_3b}{\texttt{https://huggingface.co/syzymon/long\_llama\_3b}}. We note that our checkpoint is backward compatible, i.e. can be used with any existing LLaMA inference code (both in Hugging Face and other implementations), albeit without long-context capabilities. 

\subsection{Experimental setup}\label{sec:llama_exp_setup}
 The architecture of the models is the same as OpenLLaMAs, see \cite{openlm2023openllama} {and Appendix~\ref{appendix:llamaarch}}. We use $\mathcal{L} = \{6, 12, 18\}$ (resp.  $\mathcal{L} = \{8, 16, 24\}$) as the memory layers for $3B$ (resp. $7B$) \textsc{LongLLaMA} model. We fine-tune the models on $10B$ (resp. $3B$) tokens using \textsc{FoT}, $8k$ context length and our dataset mixture based on RedPajama \citep{together2023redpajama}, see Appendix \ref{appendix:llamaDataset}.

 There are three minor differences from the standard \textsc{FoT} procedure. First, we retain the positional encodings in the local context of the memory layers (this is not necessary for \textsc{FoT}, but makes our checkpoints fully compatible with any existing LLaMA inference codebase). To be more precise, queries and keys from the local context (up to $2K$ tokens) receive the standard LLaMA rotary positional encoding, whereas memory keys are encoded as if they had position 0 in the local context window. Second, we use dense attention instead of the kNN retrieval, as we found only marginal performance differences, and it is simpler to implement. Third, we modify the crossbatch training procedure to have more fine-grained control over the number of additional contexts and the ratio of positive to negative samples. All these differences are detailed in Appendix \ref{appendix:methodsopenllama}.

\subsection{Context length extrapolation on the passkey retrieval task}

We first measure the effective context length of \textsc{LongLLaMA}, namely the distance for which tokens can effectively attend each other. We use passkey retrieval introduced in \citep{mohtashami2023landmark}, a synthetic task designed to measure this property. In this task, the model has to retrieve a passkey placed randomly in a long prompt. Results are shown in Figure \ref{fig:passkey_100k} - importantly, our $3B$ model is capable of solving this task much beyond its training context length $8K$, achieving $94.5\%$ accuracy for prompts of length $100k$ and $73\%$ for $256k$. 

\subsection{Question answering over research papers} In Table \ref{tab:quasper_results} we present the performance  on the validation set of Qasper \citep{dasigi2021dataset} from SCROLLS \citep{shaham2022scrolls} and compare our results to LongChat 7B \citep{longchat2023} and two baseline short-context models. We note that our model shows gains from increased context length.

\subsection{Improving few-shot learning accuracy with longer context} \label{sec:llama_fewshot}
 We measure long-context capabilities of these models on two downstream tasks, TREC question classification \citep{li-roth-2002-learning, hovy-etal-2001-toward} and WebQS question answering \citep{berant-etal-2013-semantic}. We follow the experimental setup of \citep{hao2022stuctured}. Namely,  we few-shot prompt the models with as many demonstration examples as possible up to the given context length. We do not use structured prompting like in \citep{hao2022stuctured} - instead, we directly provide all demonstrations in context. 

We observe significant accuracy gains from longer contexts on TREC and some improvements on WebQS (see Table \ref{tab:trec_web_qs_improve}). The TREC dataset consists of $50$ classes. A model is tasked to predict the class label given in-context examples. Only $100$ examples fit the standard context length ($2K$); it is not unusual that no class example is present for a given question, making the task impossible. Increasing the context length and the number of examples mitigates this risk. Moreover, having more demonstrations of the given class is also likely to be beneficial.

\subsection{Comparison to standard long-context fine-tuning}

In this section, we compare \textsc{FoT} to standard long-context fine-tuning, showing that it already achieves better performance for the context length used for fine-tuning and, importantly, that it can extrapolate beyond this context length, which is not the case for the baseline.

For comparisons, we fine-tune two models, one trained with \textsc{FoT} and another one (baseline) with standard fine-tuning (done similarly to \citep{MosaicML2023Introducing, XGen}). In both cases, we use $3B$ models fine-tuned on $1B$ tokens using the $4K$ context length. We evaluate both models on a number of few-shot downstream tasks in the setting described in Section \ref{sec:llama_fewshot}. 

In most cases, see Table \ref{tab:fot_vs_std_ft}, we observe accuracy improvements when more few-shot demonstrations are provided in the extended context (from $2K$ used by OpenLLaMA to $4K$ used in our fine-tuning). On TREC, the gains from additional context are significant for both models, while on WebQS, the standard fine-tuning baseline does not provide any improvement from extended context. Notably, the model fine-tuned with \textsc{FoT} enjoys further accuracy gains when evaluated with context lengths beyond its training length ($6K$ and $8K$). This shows extrapolation capabilities of \textsc{FoT}, which are not present in the baseline (see e.g. Figure \ref{fig:passkey_100k}).

\subsection{Performance on short-context tasks} Fine-tuning for longer contexts could hurt performance on the original context length ($2K$), as the training data distribution changes. We show that this is not the case for the \textsc{LongLLaMA} models by evaluating them using the LM Evaluation Harness library \citep{eval-harness}. On most tasks, the performance is kept intact; see Appendix \ref{appendix:lm_eval_harness} for details. This also confirms that \textsc{LongLLaMA}s could be used as a drop-in replacement of LLaMA models as they are compatible with the original LLaMA inference code.

\section{ Analysis of \textsc{FoT}}
In this section, we perform extensive experiments on smaller models to analyze and further validate our approach. In particular, we answer the following questions: (1) How does \textsc{FoT} perform when scaling the context length at inference time? (2) Can \textsc{FoT} be used to extend the context length of an existing, pre-trained model? (3) How effectively can it handle distractions, and how does this capability translate to enhanced performance in long-context language modeling tasks?  Moreover, we provide ablation studies of our method and additional analysis.

\subsection{Experimental setup}\label{sec:exp_setup}
\textbf{Architecture} For experiments described in this section we use decoder-only Transformer \citep{vaswani2017attention} models with $12$ layers and $184M$ parameters (unless stated otherwise). Following \cite{wu2022memorizing}; we pick $\ell=8$ as the memory attention layer. We tune $k=128$, the number of top keys retrieved by kNN. In most experiments, we start training with a small crossbatch dimension $d\leq 8$ and switch to $d\geq 64$ after some training. For more details about the architecture and hyperparameters, see Appendix \ref{appendix:architecture} and Appendix \ref{appendix:hyperparamters}. 

\textbf{Evaluation}~We distinguish two evaluation settings: single-document (abbreviated to single-doc) and multi-document (abbreviated to multi-doc). The single-doc setting is typically used for evaluating models that process long contexts. Here, we clear the memory for each new document, ensuring that only the current document is available in the context. The multi-doc setting retains memory across multiple documents without resets. This scenario {tests whether the model can ignore irrelevant information and focus on the relevant data, which can be useful in setups like repository-level code generation.}

\textbf{Datasets}~We evaluate on the following long-context language modeling datasets: PG-19 (English books), arXiv (mathematical papers), GitHub (code), and Isabelle (formal proofs). PG-19~\citep{raecompressive2019} is a large dataset of English-language books published prior to 1919, sourced from the Project Gutenberg archive. This dataset is a well-established benchmark for evaluating long-context language models~\citep{sun2021longrange}. The arXiv dataset contains \LaTeX{} source of papers labeled as "Mathematics" that were obtained by downloading articles through the arXiv Bulk Data Access. The token count per paper in this dataset is comparable to that of a book in PG19. For details on the remaining datasets, refer to Appendix \ref{appendix:datasets}.

\subsection{\textsc{FoT} fine-tuning and context length extrapolation}\label{sec:fft}
\textsc{FoT} is a minimal modification to the standard transformer architecture; therefore, it is possible to fine-tune existing models to endow them with a longer context length via the memory attention layer, as we already demonstrated in Section \ref{sec:longLama}. In this section, we deepen this analysis (on a smaller model) by studying perplexity improvements on various datasets.

As a base model, we use a standard transformer model pre-trained for $100k$ steps with context of $1K$ tokens using the standard objective and fine-tune with the \textsc{FoT} objective (i.e. crossbatch). The data used for both fine-tuning and pre-training is the C4 dataset \cite{c42019t5} (we omit documents shorter than $2K$ tokens). The fine-tuning phase takes $10k$ steps. We use the crossbatch dimension $d=128$ and local context of $1K$ tokens (context is $2K$ during training). We evaluate models in a \emph{zero-shot} way on $4$ language modeling datasets, which require long context: arXiv, PG-19, GitHub and Isabelle, see Section \ref{sec:exp_setup} and Appendix \ref{appendix:hyperparamters} for details. 

In Table \ref{tab:fft1}, we observe that \textsc{FoT} enjoys steady perplexity gains up to $64K$ tokens, although it was fine-tuned only with the $2K$ total differentiable context length. We compare the model perplexity to the following baselines: Memorizing Transformer (MT) \citep{wu2022memorizing} fine-tuned with the local context of $1K$ and memory size of $16K$, and Transformer-XL \citep{dai2019transformerxl} fine-tuned with both local context and window length of $1K$. To ensure a fair comparison, all three models are fine-tuned from the same base checkpoint. When evaluated with a context of $2K$, our method achieves results on par with the Transformer-XL baseline, which has access to the previous context in all layers, unlike MT and \textsc{FoT}. Compared to the MT baseline, we achieve better scaling when evaluated with $64K$ context length and significantly better perplexity values. Unlike MT, our method does not require training on long sequences, which is reflected by the lower perplexities of \textsc{FoT} when evaluated in the zero-shot setting. For more details, see Appendix \ref{appendix:fine_tuning}.   

{We also confirm the context extrapolation abilities using a synthetic dictionary lookup task. In this task, the model is first provided with $k_i: v_i$ mappings and then asked what value is associated with a particular key.  We train $37$M parameter  models using documents of length $512$. Figure \ref{fig:synthetic_cb_vs_baseline} shows that \textsc{FoT}, after $5$k steps of training, can effectively utilize memory consisting of $16$M tokens achieving accuracy above $92\%$. Details can be found in Appendix \ref{appendix:synthetic_task}.}

\subsection{Handling distractions in language modeling tasks}\label{sec:distractions}
 In this section, we measure how handling distractions in the multi-document setting helps in language modeling. We pick the PG-19 dataset \citep{raecompressive2019} and measure the perplexity of the next token prediction (language modeling task) when varying the size of multi-doc memory (in this case consisting of books). Intuitively, the memory tokens corresponding to the current book might be beneficial (which is also confirmed in \citep{wu2022memorizing}), while the ones from the other books are unlikely to be useful and thus are distractions. 

We observe, see Figure \ref{fig:fot_vs_baseline_multi}, that higher values of the crossbatch dimension $d$ lead to better perplexity. This aligns with the observations in Section \ref{sec:distraction_issue}, indicating that by mitigating the distraction issue, we experience benefits in language modeling.

Moreover, all versions of \textsc{FoT} are able to utilize memory and achieve much better perplexity than the standard Transformer (no memory). Unsurprisingly, perplexity increases with memory size, but we stress that this happens gracefully. In the standard variant of \textsc{FoT} (bold line), the perplexity increases only by $0.18$ when scaling to $>500k$ tokens. Importantly, the perplexity of \textsc{FoT} is close to this of Memorizing Transformer with the single-doc memory, which we treat as a soft lower bound since it is not exposed to distractions from unrelated books.

\subsection{Context length extrapolation in single-doc}\label{sec:extrapolation}

The {original} motivation behind \textsc{FoT} is to improve the multi-doc setting performance by handling distractions. Interestingly, our method also helps to extrapolate to longer contexts, even when evaluated in the single-doc setting.

To study this, we perform FoT fine-tuning (as in Section \ref{sec:fft}) and evaluate the perplexity of the resulting model on the PG-19 dataset with different context lengths in the zero-shot fashion. To deepen the analysis, we introduce an additional parameter $w$ (the number of previous contexts used in cross batch training procedure). We provide results for $w=1$ (the standard  setting for \textsc{FoT}, that corresponds to the total differentiable context being $2\cdot1024$)
and $w=2$ (corresponding to the total differentiable context  $3\cdot1024$).

We observe, see Figure \ref{fig:c4_to_pg_wd_zero_shot}, improvements when context grows, even far beyond the training context length, which reaffirms the hypothesis that \textsc{FoT} helps with extrapolation to longer contexts. Moreover, $d=2$ is significantly better than $d=1$. When comparing $d=1$ and $w=2$ to $d=2$ and $w=1$, we observe that the former is slightly better. This is natural, as the former has longer training context.

\subsection{Ablations and design choices}\label{sec:ablations}

\section{Limitations and future work}

 Our research opens a few avenues for future work. We list them as well as challenges and limitations.

\textbf{Scaling up {context}} This is by far the most important future research direction. The challenges start from purely engineering, storing more than $16$M $(key, value)$ pairs will require a distributed multi-node system. In our experiments, we use the exact kNN search, which is not scalable to large memory. Using approximate kNN search will require a lot of engineering effort, as well as careful evaluation of the impact of the approximation on the model performance. 

\textbf{Scaling up crossbatch} We observed that increasing $d$ is beneficial. In our experiments, we used $d=64$ or $d=128$, which is the maximum value that fits into the memory of a single TPUv3/TPUv2 machine, see also Appendix \ref{appendix:hardware}.  
In future work, we want to further increase $d$ as well as test on devices with bigger memory or utilize multi-node training. {We also note that crossbatch increases the training cost, but only in a subset of layers.}

\textbf{Exploring contrastive learning} The \textsc{FoT} training is inspired by rather basic contrastive learning (CL) techniques. We show that this improves the key structure so that the distraction issue is mitigated. We expect that other CL methods could be beneficial, for example, hard negative mining to utilize a larger memory during training (see \citep{lindgren2021efficient}). We leave this for future work.

\textbf{Combining with other methods} Developing long-context methods is an active research field, see Section \ref{sec:related_work}. We believe that some of these methods could be combined with \textsc{FoT}, resulting in mutually beneficial interactions.

\begin{minipage}{\linewidth}
\begin{lstlisting}[language=Python,basicstyle=\small\ttfamily , caption={Possible implementation of cross-batch. To simplify the code we assume that each document occupies two consecutive elements of the batch. A more detailed version is in Appendix \ref{appendix:code}. }, label=listing:example_cb_implemtation]

# keys from other contexts will be encoded as if they # were at the beginning of the local context pkey_fst = pos_encode_as_first(xk=key)

# local context keys encoded in the standard way pquery, pkey = pos_encode(xq=query, xk=key)

# for each element of the batch we calculate indices of # the batch that will be used in cross-batch cross_batch_rel_ids = jnp.arange(0, -num_attentions, -1) .reshape(1, -1) batch_ids = jnp.arange(0, batch_size).reshape(-1, 1) cross_batch_selector = cross_batch_rel_ids + batch_ids

# here we want other contexts cross_batch_keys = pkey_fst[cross_batch_selector[:, 1:]]

# here we concatenate local context with other contexts attention_keys = jnp.concatenate([pkey[:, None], cross_batch_keys], axis=1)

cb_attn_weights = jnp.einsum("bqhd,bckhd->bhqck", pquery, attention_keys, precision=precision)
\end{lstlisting}
\end{minipage}

\begin{ack}

\end{ack}

\end{document}
